% Reconstruction includes the finite-free coefficient-ring correction
% accepted by all four proof auditors before restoration.
\begin{proposition}[Proved: Curvewise conductors and rank-linear cohomology]
\label{D:surface-curvewise-conductor}
Let $S$ be a normal proper surface over an algebraically closed field
of characteristic $p>0$, let $Y\subsetneq S$ be closed, and let
$Q=Q(S,Y)$ be the quotient requiring tameness only at the original
codimension-one components of $Y$. There are a smooth projective
surface $\widetilde S$, an SNC divisor $B$, and a fixed effective
integral divisor $R$ supported on $B$, with the following properties.
\begin{enumerate}
\item Pullback to $V=\widetilde S\setminus B$ gives a surjection
$\pi_1(V)\to Q$. Every finite-free local system over a finite local
ring of residue characteristic different from $p$ whose monodromy
factors through $Q$ has logarithmic conductors bounded curvewise by
$R$. In particular, for every morphism $g:C\to\widetilde S$ from a
smooth curve not contained in $B$, its highest logarithmic slope at
$x\in C$ is at most $\operatorname{mult}_x(g^*R)$.
\item There is a constant $c_{S,Y}$, independent of the finite field
$\Lambda$ of characteristic $\ell\ne p$ and of the rank, such that
\[
 \dim_\Lambda H^1(Q,M)\le c_{S,Y}\dim_\Lambda M
\]
for every finite-dimensional continuous $\Lambda[Q]$-module $M$.
For an open subgroup $H\subset Q$ the bound is
$c_{S,Y}[Q:H]\dim M$.
\item These assertions and the available wild Lefschetz theorem do
not by themselves give one finitely generated group surjecting onto
$Q$: the latter theorem's finite-coefficient choices depend on the
rank and on the finite coefficient group.
\end{enumerate}
\end{proposition}
\begin{proof}
Choose a projective resolution and further blowups so that the strict
transform of the original divisorial boundary, all exceptional curves,
and the inverse images of the remaining boundary points have SNC
support $B$. We remove every exceptional curve, including those above
retained singularities. Thus $V$ maps isomorphically to a dense smooth
open of $S\setminus Y$. A connected finite etale cover of the normal
scheme $S\setminus Y$ remains connected on this dense open, which
proves the quotient assertion.

Corollary~\ref{D:surface-uniform-conductor}, applied with its
original-boundary different estimate, bounds the generic logarithmic
break at every component of $B$ by a number depending only on the
fixed pair and resolution. It is a bound for all finite quotients,
not merely for one rank. A representation of such a quotient has no
larger slope. Choose integers dominating these finitely many numbers
and let $R$ be their divisor. For smooth SNC pairs the generic
logarithmic divisor controls pullback to every curve
\citep[Theorem 3.5, Remark 3.9 and Proposition 5.8(3)]{C_HuTeyssierRam2025}.
This gives the first assertion. The finite-local-ring statement here
uses locally constant sheaves with finite free stalks, exactly the
authors' $\operatorname{Loc}_{\rm tf}$ convention; we do not silently
replace it by arbitrary finite modules. The logarithmic conductor is
the highest slope, rather than the rank-weighted Swan conductor.

For finite-field coefficients, the uniform Betti theorem on a smooth
proper compactification bounds $h^1(V,M)$ by a fixed polynomial in
the largest generic logarithmic conductor, multiplied by the rank
\citep[Theorems 7.15 and 7.34, including the proof of 7.34]{C_HuTeyssierBetti2025}.
All its conductor arguments are now fixed, so this is $c_{S,Y}\dim M$.
The continuous group $H^1$ equals etale $H^1$ for a local system,
and inflation along $\pi_1(V)\twoheadrightarrow Q$ is injective.
This proves the second assertion. Shapiro's lemma gives its open
subgroup form by multiplying the induced rank by $[Q:H]$.

Finally we retain the exact coefficient quantifiers in the corrected
wild Lefschetz result: the final version uses finite coefficients,
and its uniform choices are for fixed rank and finite coefficient
group, with dependence through $\operatorname{GL}_r(\Lambda)$
\citep[Theorems 1.4--1.5; version 3]{C_HuTeyssierLefschetz2026}.
No assertion uniform over all these groups follows. The explicit
bounded-break simple-product construction in
Proposition~\ref{D:bounded-ramification-generation-obstruction}
also prevents inferring full finite generation from these conductor
and linear-cohomology bounds alone.
\end{proof}
